\documentclass[12pt]{article}
\usepackage[utf8]{inputenc}
\usepackage[T1]{fontenc}
\usepackage{hyperref}
\title{Rwa Language Resources for K-12 Curriculum Development}
\author{Compiled by Claude Sonnet 5.5 (claude-sonnet-5-5) for Andr\'as Kornai}
\date{}
\begin{document}
\maketitle

\section{Introduction}
Rwa (ISO 639-3 \texttt{rwk}) is a Bantu language of northern Tanzania, spoken at the foot of Mount Meru. It is also called Kirwa, Kirwo, Rwo, Meru, Kimeru or Mero. A Bantu-language database places it in the Chaga group (Bantu E60) and estimates about 200,000 speakers. Other sources give about 100,000. Two Localazy pages describe a ``Rwa'' in Rwanda and Burundi, which appears to be a mislabelled entry for another language, and they are not used here.

\section{Primary Educational Resources}
No primers, textbooks or school programmes were found.

\section{Linguistic and Reference Materials}
\begin{itemize}
\item Bantu DARC entry, ``Rwa (E621A)'': a classification and basic data page. \url{https://bantudarc.aa-ken.jp/rwa.html}
\item Glottolog reference record: Rubanza, Yunus Ismail 2008. According to the search summary, it is a Kimeru lexicon from the University of Dar es Salaam catalogued under \texttt{rwk}. Its contents were not checked. \url{https://glottolog.org/resource/reference/id/560504}
\end{itemize}

\section{Religious and Community Resources}
\begin{itemize}
\item Scripture Earth page for \texttt{rwk}: it lists the alternative language names and may point to scripture resources. \url{https://www.scriptureearth.org/00eng.php?iso=rwk}
\end{itemize}

\section{Digital and Audio Resources}
\begin{itemize}
\item ICU/CLDR locale entry ``Rwa (Tanzania)'' (\texttt{rwk\_TZ}): software locale data. \url{https://www.localeplanet.com/icu/rwk-TZ/index.html}
\end{itemize}

\section{Curriculum Development Context}
The language is not known to be taught in schools. Kiswahili is the medium of instruction in Tanzanian primary schools and English in secondary schools. The most promising leads are the Dar es Salaam lexicon and any scripture materials. Both should be checked directly before relying on them.
\end{document}